%-------------------------
% Plantilla de CV de una columna para que el texto se extraiga limpio.
%
% Basada en «Jake's Resume» (https://github.com/jakegut/resume),
% Copyright (c) 2020 Jake Gutierrez, licencia MIT, que a su vez se basa en
% sb2nov/resume (https://github.com/sb2nov/resume), Copyright (c) Sourabh Bajaj,
% licencia MIT. Los avisos completos están en LICENSE.
%
% Cambios: A4, español, sin iconos, mapa Unicode activado, sin guiones de
% corte, viñetas y fechas que se extraen en orden.
%-------------------------
\documentclass[a4paper,10pt]{article}

\usepackage[T1]{fontenc}
\usepackage[utf8]{inputenc}
\usepackage{lmodern}
\usepackage{latexsym}
\usepackage[empty]{fullpage}
\usepackage{titlesec}
\usepackage[usenames,dvipsnames]{color}
\usepackage{enumitem}
\usepackage[hidelinks]{hyperref}
\usepackage{fancyhdr}
\usepackage[spanish,es-noshorthands,es-nolayout]{babel}
\usepackage{tabularx}

% Mapa Unicode: sin esto, «fi» o «ñ» pueden salir como glifos raros al extraer el texto
\input{glyphtounicode}
\pdfgentounicode=1

\pagestyle{fancy}
\fancyhf{}
\fancyfoot{}
\renewcommand{\headrulewidth}{0pt}
\renewcommand{\footrulewidth}{0pt}

% Márgenes
\addtolength{\oddsidemargin}{-0.55in}
\addtolength{\evensidemargin}{-0.55in}
\addtolength{\textwidth}{1.1in}
\addtolength{\topmargin}{-.65in}
\addtolength{\textheight}{1.35in}

\urlstyle{same}
\raggedbottom
\raggedright
\setlength{\tabcolsep}{0in}

% Sin guiones de corte: ensucian el texto que lee un sistema automático
\hyphenpenalty=10000
\exhyphenpenalty=10000

% Títulos de sección
\titleformat{\section}{
  \vspace{-4pt}\scshape\raggedright\large
}{}{0em}{}[\color{black}\titlerule \vspace{-5pt}]

% Comandos
\newcommand{\resumeItem}[1]{\item\small{{#1 \vspace{-2pt}}}}

% Entrada con cuatro datos: {centro o empresa}{lugar}{título o puesto}{fechas}
% El original pone lugar y fechas en una tabla de dos filas a la derecha, y al
% extraer el texto salían detrás de los puntos. Aquí el lugar va junto al centro
% y solo la fecha va a la derecha, así se extrae justo después de su entrada.
\newcommand{\resumeSubheading}[4]{
  \vspace{-2pt}\item
    \textbf{#1}, #2 \hfill \textit{\small #4} \\
    \textit{\small #3}\vspace{-7pt}
}

% Proyecto: {nombre | tecnologías}{fechas}
\newcommand{\resumeProjectHeading}[2]{
    \item
    \begin{tabular*}{0.97\textwidth}{l@{\extracolsep{\fill}}r}
      \small#1 & \small #2 \\
    \end{tabular*}\vspace{-7pt}
}

% Viñetas que se extraen en su sitio
\renewcommand\labelitemi{\textbullet}
\renewcommand\labelitemii{$\vcenter{\hbox{\tiny$\bullet$}}$}

\newcommand{\resumeSubHeadingListStart}{\begin{itemize}[leftmargin=0.15in, label={}]}
\newcommand{\resumeSubHeadingListEnd}{\end{itemize}}
\newcommand{\resumeItemListStart}{\begin{itemize}}
\newcommand{\resumeItemListEnd}{\end{itemize}\vspace{-5pt}}
